% Source: https://artofproblemsolving.com/community/c3h1582862_17_geometry_questions
% Stable identity: aops-topic-1582862-q13
In $15$ minutes she completes $\dfrac{15}{90}=\dfrac{1}{6}$ of an orbit, so the point of the equator directly below her turns through $\dfrac{1}{6}\cdot360^\circ=60^\circ$ about Earth's centre. She passes over $\dfrac{1}{6}$ of the equator:

\begin{asy}
usepackage("amsmath");
size(4.5cm);
defaultpen(fontsize(10pt));
pair O=(0,0);real R=2;
draw(circle(O,R),linewidth(0.9));draw(circle(O,2.5),dashed);
draw(O--R*dir(90));draw(O--R*dir(30));
draw(arc(O,R,30,90),red+linewidth(2));
draw(arc(O,0.5,30,90));label("$60^\circ$",0.85*dir(60));
label("$8000$ mi",(-R,0)--(R,0),S);draw((-R,0)--(R,0),dotted);
label("orbit",2.5*dir(215),dir(215),fontsize(8pt));
label("\mbox{}",(min(currentpicture,user=true)+max(currentpicture,user=true))/2,black); // keeps the text after the figure black when LaTeX typesets the labels inline
\end{asy}

\[
\dfrac{1}{6}\cdot\pi\cdot8000=\dfrac{4000\pi}{3}\approx\boxed{4189\text{ miles}}.
\]
(Her height above Earth does not matter: the arc on the ground below her is determined by the angle she turns through.)
